\documentclass[../../../../main.tex]{subfiles}
\begin{document}

図で，$g$ は水平面に対する傾き $\tan\alpha$ が $1/2$ であるような定直線とし，
$OA$，$AB$ は $A$ で(ちょうつがいで)連結された長さの等しい棒で，
その端 $O$ は $g$ 上の定点に固定され，
$OA$ は $g$ を含む鉛直面内で自由に回転し，
他の端 $B$ は $g$ 上を動くことができるようになっている．

このとき，折れ線 $OAB$ の重心 $G$($OA$，$AB$ の中点を結ぶ線分の中点)が最低になるのは，
$OA$ の水平面となす傾き $\tan\theta$ がいくらになるときか．

\begin{figure}[htbp]
  \centering
  \begin{tikzpicture}[scale=1.5, >=stealth]
    \draw[thick] (-1.5, 0) -- (2.5, 0);

    \draw[thick] (-1.5, -0.75) -- (2.5, 1.25) node[above right] {$g$};
    \fill (1.5, 0.75) circle (1.5pt) node[above left] {$O$};

    \draw (1.5, 0.75) +(0:0.5) arc (0:26.57:0.5);
    \node at ($(1.5, 0.75)+(13.28:0.7)$) {$\alpha$};

    \coordinate (O) at (1.5, 0.75);
    \coordinate (A) at (1.0, -0.8);
    \coordinate (B) at (-1.0, -0.5);

    \draw[line width=2pt] (O) -- (A) -- (B);

    \fill[white] (A) circle (3pt);
    \draw (A) circle (3pt);
    \node[below right] at (A) {$A$};

    \node[below left] at (B) {$B$};

    \coordinate (M1) at ($(O)!0.5!(A)$);
    \coordinate (M2) at ($(A)!0.5!(B)$);
    \coordinate (G) at ($(M1)!0.5!(M2)$);

    \fill (G) circle (1.5pt) node[left=2pt] {$G$};

    \draw (O) +(-153.43:0.6) arc (-153.43:-108.43:0.6);
    \node at ($(O)+(-130:0.8)$) {$\theta$};
  \end{tikzpicture}
\end{figure}

\end{document}
